\documentclass{article}
\usepackage{amsmath,amssymb}
\title{Relative pure-endpoint laws for native synchronization}
\date{15 November 2026}
\begin{document}\maketitle
\paragraph{Status.} Independently reviewed application of the existing native cone-creation and preservation criterion. The purpose is to verify the original relative pure interface, rather than leave synchronization as unrelated fixed-interface functors. No novelty is claimed for change-of-endpoint whiskering or preservation over an identity base.

\paragraph{Data.} Let $\mathcal S$ be the retained state bicategory and $\mathcal K$ the process bicategory. For each external pure interface $A$ and $B$ let $P_A,Q_B:\mathcal S\to\mathcal K$ be the specified pure-action pseudofunctors. An original pure map $a:A'\to A$ induces a pseudonatural transformation $p_a:P_{A'}\Rightarrow P_A$; a pure map $b:B\to B'$ induces $q_b:Q_B\Rightarrow Q_{B'}$. Their component arrows need not be equivalences. Their pseudonaturality cells and the coherent identity/composition modifications are those of the given pure actions, not new free choices. Assume the native dialgebra homs $D(A,B)$ have the selected incoming lax bilimits created over the same selected state cones; in particular the corresponding $Q_B$ pseudofunctors preserve those cones as required by the creation theorem.

\paragraph{Actual endpoint transport.} Define
\[
 T_{a,b}:D(A,B)\longrightarrow D(A',B'),\qquad
 (X,f)\longmapsto(X,q_b(X)f p_a(X)).
\tag{1}
\]
For $(h,\alpha)$, keep the exact state map $h$. The target native square is
\[
 Q_{B'}h\,q_b(X)f p_a(X)
 \cong q_b(Y)Q_Bh\,f p_a(X)
 \xRightarrow{q_b(Y)\alpha p_a(X)}
 q_b(Y)gP_Ah\,p_a(X)
 \cong q_b(Y)g p_a(Y)P_{A'}h.
\tag{2}
\]
The outer comparisons are pseudonaturality cells of $q_b$ and $p_a$. No component of either transformation is inverted. On a native two-cell retain its state two-cell; its new native equation follows by whiskering the original equation and using pseudonaturality on that two-cell. Composition and identity comparisons have the original state-base constraints and the action pseudonaturality comparisons. Their equations follow by the composition laws of $p_a,q_b$ and bicategory coherence. Thus this is a native pseudofunctor over the identity state base, even when the pure realization is nonfaithful.

\paragraph{Selected cone preservation.} Transport the native universal cone of a diagram $F:J\to D(A,B)$ by (1). Its underlying state cone is unchanged, since every state object, map and two-cell is retained. Its native projection squares remain invertible: (2) is a paste of the original invertible projection square and invertible pseudonaturality cells. The native creation theorem in $D(A',B')$ therefore recognizes it as a universal cone. Equivalently this is the earlier preservation criterion with base functor $1_{\mathcal S}$, which preserves the selected cones, and the just-proved pseudo-arrow preservation. No separate preservation hypothesis on the component arrow $q_b(X)$ is required.

There is consequently a cone-compatible pseudonatural equivalence
\[
 T_{a,b}\,L_J^{A,B}\simeq
 L_J^{A',B'}\,(T_{a,b})_*.
\tag{3}
\]
It is constructed by the actual cone equivalences, not by asserting equality of their chosen representing state objects. Naturality holds on all native diagram arrows and modifications. It extends to directed lax-body maps and reverse shape changes because both routes send their cone modifications by precisely (2) and the same original directed cone cells. The component maps of $a,b$ may remain noninvertible throughout.

\paragraph{Coherent relative action.} For composable pure maps, the original action modifications give the specified pseudonatural comparison between composite transports $T$ and transport by the composite endpoint maps. Under (3), its two synchronization pastes have the same projections and the same native action comparison at every diagram object. Full faithfulness yields their invertible comparison modification, and proves its unit and associativity equations. This lifts the actual original relative action with its frames; it does not add new shape maps to the external pure category or infer state-map equality from endpoint images.

The directed sequential evaluation comparison is compatible with these relative actions. For composable input diagrams, form the two routes using their appropriate external endpoint transports and the original pure tightening/centrality comparison at the shared interface. Project the result to each product-shape pair. The two routes become the same original tightening/action equation on the actual pair of process bodies and their native squares. Once that original comparison has been admitted, the cone equivalence produces the comparison modification and detects its equations. The same statement applies to an admitted parallel operation with its specified pure-action compatibility. No endpoint comparison or middle-interface sliding is postulated if the original action package does not supply it.

\paragraph{Connection to varying-state feedback.} The separate native-dialgebra diagram presentation identifies a varying retained-state diagram as $X:J\to\mathcal S$ together with a directed lax body $P_AX\Rightarrow Q_BX$. Formula (1) transports exactly that body along the original pure external maps while retaining every $X$ and every native witness. Equation (3) therefore verifies genuine indexed process semantics, not just a transformation of labels attached to unrelated hom categories.

Canonical finite state joining, its normalization and join-then-synchronize evaluation use the already specified relative pure-action comparisons. Combining those comparisons with (3) gives the external pure-endpoint law for join-first evaluation. This does not add a generic semantic feedback evaluator: its action on arbitrary state-uniformity squares still must be the actual admitted native functor. Nor does it prove strong synchronization compatibility for a varying state tensor. It establishes the original external pure interface and its directed assembly compatibility under the stated cone-creation hypotheses.
\end{document}
